%=====================================================

% lista de símbolos

\begin{listasimb}

\begin{longtable}[l]{p{0.22\linewidth}p{0.70\linewidth}}
$a$, $b$ & Dimensões laterais da cavidade retangular formada pelos planos\\
$c_0$ & Velocidade da luz no vácuo\\
$C_{\text{plate}}$ & Capacitância de placas paralelas da cavidade\\
$f$ & Frequência\\
$f_{mn}$ & Frequência do modo ressonante de índices $m,n$ da cavidade\\
$h$ & Altura da cavidade (espessura do dielétrico)\\
$\mathbf{L}$ & Função de perda (\textit{loss}) do modelo\\
$\lambda_k$ & Peso do $k$-ésimo termo da função de perda composta\\
$N$ & Número de amostras do conjunto de dados\\
$n_f$ & Número de pontos de frequência\\
$R^2$ & Coeficiente de determinação\\
$\mathbf{S}(f)$ & Matriz de parâmetros de espalhamento\\
$\sigma$ & Condutividade elétrica do condutor\\
$\tan\delta$ & Tangente de perdas do dielétrico\\
$\varepsilon_0$ & Permissividade elétrica do vácuo\\
$\varepsilon_r$ & Permissividade relativa do dielétrico\\
$\mathbf{x}$ & Vetor de parâmetros de projeto (entrada do modelo)\\
$\mathbf{Z}(f)$ & Matriz de impedância em função da frequência\\
$Z_{11}(f)$ & Autoimpedância vista na porta 1\\
$|Z|_{\max}$ & Máximo do módulo da impedância na faixa de análise\\
$Z_{\text{target}}$ & Impedância-alvo de projeto da PDN\\
$\omega$ & Frequência angular, $\omega = 2\pi f$\\
\end{longtable}

\end{listasimb}

%=====================================================
